% BEGIN VOL08-EXPANSION VIII14
\section{Supplementary worked examples}
\label{sec:viii14-pedagogy}

\begin{example}[Rose graph]\label{ex:viii14-ped-01}
The wedge $R_r=\bigvee_{i=1}^r S^1$ has one vertex and $r$ oriented loop edges.  Reduced edge words give its loop classes, yielding $\pi_1(R_r)\cong F_r$.
\end{example}

\begin{example}[Universal covering tree]\label{ex:viii14-ped-02}
The universal cover of a connected graph is a tree.  For the rose $R_r$, vertices can be labeled by reduced words in $F_r$, and an edge appends one generator before free reduction.
\end{example}

\begin{example}[Subgroups as covering graphs]\label{ex:viii14-ped-03}
Under covering-space classification, a subgroup $H\le F_r$ determines a connected covering graph of the rose.  Cosets describe the fiber, and loops based at a chosen lift encode exactly the subgroup $H$.
\end{example}

\section{Graded supplementary exercises}
The following exercises progress from direct computation and construction to proof, hypothesis testing, application, and synthesis.

\subsection*{Standard computations and constructions}

\begin{exercise}[Rank-one rose]\label{exr:viii14-ped-01}
What is $\pi_1(R_1)$?
\end{exercise}
\begin{hint}
$R_1=S^1$.
\end{hint}
\begin{solution}
$\pi_1(R_1)\cong\mathbb Z=F_1$.
\end{solution}

\begin{exercise}[Rank-two rose]\label{exr:viii14-ped-02}
Name the group $\pi_1(R_2)$.
\end{exercise}
\begin{hint}
There are two loop generators and no 2-cell relations.
\end{hint}
\begin{solution}
It is the free group $F_2=\langle a,b\rangle$.
\end{solution}

\begin{exercise}[Tree group]\label{exr:viii14-ped-03}
Compute the fundamental group of a tree.
\end{exercise}
\begin{hint}
A tree contracts to a point.
\end{hint}
\begin{solution}
It is trivial.
\end{solution}

\begin{exercise}[Index and sheets]\label{exr:viii14-ped-04}
If $H\le F_r$ has finite index $d$, how many sheets does the associated connected cover have?
\end{exercise}
\begin{hint}
Fibers correspond to cosets.
\end{hint}
\begin{solution}
It is a $d$-sheeted covering.
\end{solution}

\begin{exercise}[Reduced word]\label{exr:viii14-ped-05}
Reduce the word $ab b^{-1}a^{-1}c$ in a free group.
\end{exercise}
\begin{hint}
Cancel adjacent inverse pairs.
\end{hint}
\begin{solution}
It reduces to $c$.
\end{solution}

\subsection*{Proofs}

\begin{exercise}[Graph universal cover is tree]\label{exr:viii14-ped-06}
Prove a simply connected connected graph has no cycles and hence is a tree.
\end{exercise}
\begin{hint}
A cycle would define a nontrivial loop after removing backtracking.
\end{hint}
\begin{solution}
Any embedded cycle contributes a nontrivial free-group element; simply connectedness forbids cycles, so a connected simply connected graph is a tree.
\end{solution}

\begin{exercise}[Rose free group]\label{exr:viii14-ped-07}
Prove $\pi_1(R_r)\cong F_r$.
\end{exercise}
\begin{hint}
Use van Kampen or reduced edge loops.
\end{hint}
\begin{solution}
Each oriented loop edge gives a generator; endpoint homotopy cancels only immediate backtracking, producing precisely reduced words with no further relations.
\end{solution}

\begin{exercise}[Subgroup from cover]\label{exr:viii14-ped-08}
Show $p_*\pi_1(E,e_0)$ is a subgroup of $\pi_1(B,b_0)$ for a based cover.
\end{exercise}
\begin{hint}
$p_*$ is a group homomorphism and injective for coverings.
\end{hint}
\begin{solution}
The image of a homomorphism is a subgroup; covering path lifting also shows $p_*$ is injective by lifting a null-homotopy.
\end{solution}

\begin{exercise}[Coset fiber]\label{exr:viii14-ped-09}
Explain why fiber points correspond to cosets of the covering subgroup.
\end{exercise}
\begin{hint}
Transport the base lift along paths representing group elements.
\end{hint}
\begin{solution}
Two loop classes send the base lift to the same endpoint exactly when their quotient lies in the subgroup of loops that lift closed; thus endpoints are indexed by the appropriate cosets.
\end{solution}

\subsection*{Counterexamples and hypothesis tests}

\begin{exercise}[All graph groups abelian]\label{exr:viii14-ped-10}
Test: fundamental groups of graphs are always abelian.
\end{exercise}
\begin{hint}
$F_2$ is not abelian.
\end{hint}
\begin{solution}
False. Connected graphs have free fundamental groups, usually nonabelian.
\end{solution}

\begin{exercise}[Universal cover finite]\label{exr:viii14-ped-11}
Test: the universal cover of a finite graph is always finite.
\end{exercise}
\begin{hint}
Consider a circle.
\end{hint}
\begin{solution}
False. The universal cover of $S^1$ is an infinite line; nontrivial free groups generally give infinite trees.
\end{solution}

\begin{exercise}[Subgroups of free groups need not free]\label{exr:viii14-ped-12}
Test: a subgroup of a free group can have torsion and need not be free.
\end{exercise}
\begin{hint}
Use covering graphs.
\end{hint}
\begin{solution}
False. Nielsen--Schreier says every subgroup of a free group is free; covering graphs provide a topological proof.
\end{solution}

\subsection*{Applications and investigations}

\begin{exercise}[Automata viewpoint]\label{exr:viii14-ped-13}
How can a finite covering graph encode a finite-index subgroup of a free group?
\end{exercise}
\begin{hint}
Edges labeled by generators define transitions between sheets/cosets.
\end{hint}
\begin{solution}
Choose a base vertex; reading a word follows labeled edges, and the accepted words returning to the base vertex are precisely the subgroup.  The finite cover acts like a deterministic finite-state machine for coset action.
\end{solution}

\begin{exercise}[Spanning tree computation]\label{exr:viii14-ped-14}
Explain how a spanning tree gives a basis for the fundamental group of a finite connected graph.
\end{exercise}
\begin{hint}
Tree edges create no cycles.
\end{hint}
\begin{solution}
Collapse the spanning tree to a point; each remaining edge becomes one loop, so the graph is homotopy equivalent to a wedge of $E-V+1$ circles and has free rank $E-V+1$.
\end{solution}

\subsection*{Challenge problems}

\begin{exercise}[Nielsen--Schreier rank]\label{exr:viii14-ped-15}
If $H$ has index $d$ in $F_r$, compute the rank predicted by a $d$-sheeted covering of the rose.
\end{exercise}
\begin{hint}
Count vertices and edges of the covering graph.
\end{hint}
\begin{solution}
The cover has $V=d$ and $E=dr$, so its free rank is $E-V+1=d(r-1)+1$.
\end{solution}

\begin{exercise}[Covering graph synthesis]\label{exr:viii14-ped-16}
Explain how lifting, deck actions, and subgroup classification meet in covering graphs.
\end{exercise}
\begin{hint}
Loops act on the fiber, closed lifts form a subgroup, symmetries normalize that subgroup.
\end{hint}
\begin{solution}
Path lifting gives the monodromy action of $F_r$ on fiber vertices; the stabilizer of the base vertex is the subgroup defining the cover; global graph automorphisms over the base form the deck group, governed by the subgroup normalizer.
\end{solution}

% END VOL08-EXPANSION VIII14
